\section{Related Work}
\label{sec:related}

\paragraph{Autonomous Research Benchmarks.}
Autonomous research benchmarks evaluate progress at different levels of an AI system. Task-level evaluations measure agents' ability to conduct ML experiments and solve AI R\&D problems~\citep{mlagentbench,mlebench,rebench}. Other evaluations focus on the supporting harness and infrastructure~\citep{harnessdev,phibench}, or on the agent's ability to learn from broad goals and interaction feedback~\citep{aspire,s3gym}. For autonomous post-training, the object of improvement is the target model: PostTrainBench permits end-to-end language-model post-training, whereas RSIBench-Data fixes the learning stack to isolate training-data research~\citep{posttrainbench,rsibenchdata}. \benchmark{} extends this setting to multimodal tasks, assessing final target and non-target performance, iterative candidate quality and submission choice, and research integrity.

%

\paragraph{Research Agent Harnesses.}
Research agent harnesses must organize experiments and make accumulated evidence useful for subsequent decisions. End-to-end workflows coordinate idea generation, implementation, and training~\citep{aiscientist,aibuildai}, while search-based systems structure the choice of experiments~\citep{aide}. As experimentation continues, preserving and reusing research state becomes equally important. Hierarchical knowledge retrieval and persistent hypothesis trees connect earlier findings to later decisions~\citep{aibuildai2,arbor}; discovery replay and reusable environment memory serve a related role in autonomous exploration~\citep{dreamrsi,rsiagent}. ModularRSI refactors Terminus-2 into behavioral modules and evolves their source functions~\citep{wu2026modularrsi}. Building on research orchestration and knowledge reuse, \harness{} focuses on improving multimodal models through post-training, connecting multimodal evidence-based diagnosis with training interventions and model selection.
